% Part: first-order-logic
% Chapter: models-theories
% Section: expressing-relations

\documentclass[../../../include/open-logic-section]{subfiles}

\begin{document}

\olfileid[pt]{fol}{mat}{exr}

\olsection{Expressando Relações em \article{structure}
  \printtoken{S}{structure}}

\begin{explain}
Um uso principal que se pode dar às !!{formula}s é expressar propriedades e
relações em !!a{structure}~$\Struct M$ em termos dos primitivos da
linguagem~$\Lang L$ de~$\Struct M$.  Com isso queremos dizer o seguinte:
o !!{domain} de $\Struct M$ é um conjunto de objetos.  Os !!{constant}s,
!!{function}s e !!{predicate}s são interpretados em~$\Struct M$ por
alguns objetos em~$\Domain M$, funções sobre~$\Domain M$, e relações
sobre~$\Domain M$.  Por exemplo, se $\Obj A^2_0$ está em $\Lang L$, então
$\Struct M$ atribui a ele uma relação~$R = \Assign{{\Obj
  A^2_0}}{M}$. Então a fórmula $\Atom{\Obj A^2_0}{\Obj v_1, \Obj v_2}$
\emph{expressa} exatamente essa relação, no seguinte sentido: se uma
atribuição de variáveis~$s$ mapeia $\Obj v_1$ para $a \in \Domain{M}$ e
$\Obj v_2$ para $b \in \Domain M$, então
\[
Rab \text{\quad sse\quad} \Sat{M}{\Atom{\Obj A^2_0}{\Obj v_1, \Obj v_2}}[s].
\]
Note que temos que envolver atribuições de variáveis aqui: não podemos simplesmente
dizer ``$Rab$ sse $\Sat{M}{\Atom{\Obj A^2_0}{a, b}}$'' porque $a$ e
  $b$ não são símbolos da nossa linguagem: eles são !!{element}s
  de~$\Domain{M}$.

Como não temos apenas !!{formula}s atômicas, mas podemos combiná-las
usando os conectivos lógicos e os quantificadores, !!{formula}s mais complexas
podem definir outras relações que não estão diretamente embutidas
em~$\Struct M$.  Estamos interessados em como fazer isso e,
especificamente, quais relações podemos definir em !!a{structure}.
\end{explain}

\begin{defn}
Seja $!A(\Obj v_1,\dots, \Obj v_n)$ !!a{formula} de $\Lang L$ na
qual apenas $\Obj v_1$,\dots, $\Obj v_n$ ocorrem livres, e seja $\Struct
M$ !!a{structure} para~$\Lang L$. $!A(\Obj v_1,\dots, \Obj v_n)$
\emph{expressa a relação}~$R \subseteq \Domain M^n$ sse
\[
Ra_1\dots a_n \text{\quad sse\quad} \Sat{M}{\Atom{!A}{\Obj
    v_1,\dots, \Obj v_n}}[s]
\]
para qualquer atribuição de variáveis~$s$ com $s(\Obj v_i) = a_i$ ($i = 1,
\dots, n$).
\end{defn}

\begin{ex}
No modelo padrão da aritmética~$\Struct N$, a !!{formula} $\Obj
v_1 < \Obj v_2 \lor \eq[\Obj v_1][\Obj v_2]$ expressa a relação $\le$
sobre~$\Nat$. A !!{formula} $\eq[\Obj v_2][\Obj v_1']$
expressa a relação sucessora, isto é, a relação $R \subseteq
\Nat^2$ onde $Rnm$ vale se $m$ é o sucessor de~$n$. A fórmula
$\eq[\Obj v_1][\Obj v_2']$ expressa a relação predecessora.  As
!!{formula}s $\lexists[\Obj v_3][(\eq/[\Obj v_3][\Obj 0] \land
  \eq[\Obj v_2][(\Obj v_1 + \Obj v_3)])]$ e $\lexists[\Obj
  v_3][\eq[(\Obj v_1 + {\Obj v_3}')][v_2]]$ ambas expressam a relação $\Obj <$.
Isso significa que o símbolo de predicado~$<$ é, na verdade,
supérfluo na linguagem da aritmética; ele pode ser definido.
\end{ex}

\begin{explain}
Essa ideia não é interessante apenas em !!{structure}s específicas, mas
geralmente sempre que usamos uma linguagem para descrever um modelo ou
modelos pretendidos, isto é, quando consideramos teorias.  Essas teorias frequentemente
contêm apenas alguns !!{predicate}s como símbolos básicos, mas no domínio que elas
são usadas para descrever muitas outras relações frequentemente desempenham um
papel importante.  Se essas outras relações podem ser sistematicamente expressas pelas
relações que interpretam os !!{predicate}s básicos da linguagem, nós
dizemos que podemos \emph{defini-las} na linguagem.
\end{explain}

\begin{prob}
Encontre !!{formula}s em $\Lang L_A$ que definam as seguintes relações:
\begin{enumerate}
\item $n$ está entre $i$ e $j$;
\item $n$ divide $m$ exatamente (isto é, $m$ é um múltiplo de $n$);
\item $n$ é um número primo (isto é, nenhum número além de $1$ e $n$ divide~$n$
  exatamente).
\end{enumerate}
\end{prob}

\begin{prob}
Suponha que a fórmula $!A(\Obj v_1, \Obj v_2)$ expressa a relação $R
\subseteq \Domain M^2$ em !!a{structure}~$\Struct M$.  Encontre fórmulas
que expressem as seguintes relações:
\begin{enumerate}
\item a inversa $R^{-1}$ de $R$;
\item o produto relativo $R \mid R$;
\end{enumerate}
Você consegue encontrar uma maneira de expressar $R^+$, o fecho transitivo de~$R$?
\end{prob}

\begin{prob}
Seja $\Lang{L}$ a linguagem contendo apenas um símbolo de predicado de 2 lugares
$<$ (sem outros !!{constant}s, !!{function}s ou !!{predicate}s---
exceto, é claro,~$\eq$). Seja $\Struct{N}$ a estrutura tal que
$\Domain{N} = \Nat$, e $\Assign{<}{N} = \Setabs{\tuple{n,m}}{n <
  m}$. Prove o seguinte:
\begin{enumerate}
\item $\{ 0 \}$ é definível em $\Struct{N}$;
\item $\{ 1 \}$ é definível em $\Struct{N}$;
\item $\{ 2 \}$ é definível em $\Struct{N}$;
\item para cada $n \in \Nat$, o conjunto $\{ n \}$ é definível em
  $\Struct{N}$;
\item todo subconjunto finito de $\Domain{N}$ é definível em
  $\Struct{N}$;
\item todo subconjunto cofinito de $\Domain{N}$ é definível em
  $\Struct{N}$ (onde $X \subseteq \Nat$ é cofinito sse
  $\Nat \setminus X$ é finito).
\end{enumerate}
\end{prob}


\end{document}
